% Page 21 — Alban en sa « légende dorée »

\begin{center}
  {\Large\itshape Alban en sa « légende dorée »}
\end{center}
\vspace{1.5em}

\noindent
\lettrine{O}{n} fit à Alain de Nougaret de très belles obsèques au cours desquelles on vanta les mérites du bon baron repentant. Au cimetière paroissial, il fut inhumé sous la plus belle croix de pierre. C'est la « Croix de la Madone » (classée monument historique) que l'on peut voir encore dans l'église Notre Dame d'Alban. Malheureusement, le sommet de ce calvaire a dû pendant des années être foulé aux pieds, et de ce fait, se trouve très abîmé. Une partie du piédroit semble aussi manquer, hélas ! Mais dans toute cette histoire où se trouve donc la fameuse pierre du diable ? Patiente, Ami – Oui, patience ! La mort entraîne le jugement. Et de ciel s'ensuit ou l'enfer. A peine Alain a-t-il exhalé son dernier souffle que le démon a bondi à l'entrée du paradis pour y cueillir son homme. « Voici le contrat dûment signé, Messire Alain. Ne nous trompons pas de chemin et en route pour l'enfer, logement gratuit et chauffage central assuré, en tout temps, hiver comme été ». Malgré sa peur affreuse, Alain a crié de toutes ses forces, appelant à l'aide. Celle qu'on n'invoqua jamais en vain – « Souvenez-vous, ô très douce Vierge marie, qu'on n'a jamais entendu dire qu'aucun de ceux qui ont eu recours à vous, imploré votre assistance, réclamé votre secours, n'ait été abandonné – O Madone d'Alban, mère très aimée.
